La investigación desarrolló una estrategia de estimación del rendimiento de cebada para 59 parcelas sin etiqueta mediante información de otras 138 parcelas observadas y un conjunto de covariables ambientales disponible para las 197 unidades. La arquitectura elegida no corresponde a una regresión aislada: reúne un componente climático Elastic Net, un experto Ridge geográfico, una compuerta de composición geológica y una corrección residual restringida. 

La imputación por media se conservó tras registrar un RMSE ligeramente menor al de la mediana: 0.468138 frente a 0.468186 toneladas por hectárea. La diferencia fue mínima y no se tradujo en una reducción del MAE.  La búsqueda escalonada de parámetros también mostró que la clasificación de candidatos cambia al aumentar el número de parcelas excluidas; al completar las 138 exclusiones, la combinación de 24 vecinos y penalización Ridge 30 resultó la mejor de las finalistas examinadas.

En materia de agrupamiento, K-means integra la arquitectura Reference26 sometida a evaluación completa. No obstante, la prueba documentada sobre 46 parcelas favoreció a DBSCAN en RMSE. Por ello no cabe defender la compuerta elegida alegando una superioridad empírica universal frente a DBSCAN. Mantener la arquitectura con K-means responde a la procedencia y verificabilidad del modelo final, mientras la variante de densidad exige una comparación global homogénea antes de modificar el sistema.

La última corrección utiliza tres parcelas análogas en el espacio robusto de 26 atributos, con un factor parcial de mezcla y una compuerta de soporte para evitar cambios excesivos. La regla fija seleccionada obtuvo un RMSE de 0.462418 toneladas por hectárea y un coeficiente de determinación predictivo de 0.702693. Una política de selección interna obtuvo RMSE de 0.462892. La magnitud de mejora fue modesta, no benefició a todas las parcelas y no eliminó los errores extremos.

 Los 138 rendimientos se consultaron repetidamente durante la exploración de familias, parámetros y correcciones. En consecuencia, una parte de la aparente ventaja puede proceder de selección adaptativa; incluso la política posteriormente anidada está condicionada a candidatos ya cribados. Tampoco se conoce el rendimiento real de las 59 parcelas objetivo, por lo que sería incorrecto atribuirles las métricas de las exclusiones de entrenamiento.

Como conclusión, el sistema aporta una reconstrucción transductiva técnicamente documentada, con decisiones verificables y criterios explícitos para excluir etiquetas de evaluación. Para una aplicación financiera o productiva fuera del reto harían falta campañas independientes, calibración de incertidumbre y análisis de estabilidad por territorio. 
